\section{Range Categories, Partial Sections and Weak-Range Operators}

\subsection{Range Categories}
Cockett et al. (\cite{COCKETT2012}) define range categories from restriction categories along the following lines.

\begin{definition} 
A \term{range category} is a restriction category with an additional operator on morphisms (the range operator $\rng{\ }$), as follows
if $f: a \morph b$ in  \catcw then
$$\rng{f}: b \morph b$$
satisfying

RR.1 For $f:a \morph b$ in \catcw $$\rst{\rng{f}} = \rng{f}.$$

RR.2 For $f:a \morph b$ in \catcw $$f \circ \rng{f} = f.$$

RR.3. If $\sequentialdiag{a}{b}{c}{f}{g}$ in \catcw then
$$\rng{f \circ \rst{g}} = \rng{f} \circ \rst{g}.$$

RR.4. If $\sequentialdiag{a}{b}{c}{f}{g}$ in \catcw then
$$\rng{(\rng{f} \circ g)} = \rng{f \circ g}.$$
\end{definition}

In this note we explore range categories that satisfy an  additional condition defined by Cockett et al. as RR.5. We will use the folowing terminology:
\newpage
\begin{definition}
An \textit{RR.5 range category} is a range category that additionally satisfies:

RR.5 if $\equaliser{a}{f}{\paralleldiag{b}{c}{g}{h}}$ then
 $$f \circ g = f \circ h \Rightarrow  \rng{f} \circ g = \rng{f} \circ h.$$
 \end{definition}

\begin{lemma}
\llabel{unirangelemma}
If $f:a \morph a$ in a range category \catcw then for any morphism 
$f: a \morph b$
\begin{enumerate}[(i)] 
\item $\rng{\rng{f}}=\rng{f}$,
\item $\rng{f} \comp \rng{f} = \rng{f}$,
\item \llabel{lem:unirange:hatbar} $\rng{\rst{f}}=\rst{f}$.
\end{enumerate}
\end{lemma}
\begin{proof}
Straightforward.
\end{proof}


\begin{lemma}\commentary{new : 3rd July 2026}
\llabel{fgJoinLemma}
 If \catcw is a category with support  and
 if $\sequentialdiag{a}{b}{c}{f}{g}$ in \catcw  then
$
\rng{f \circ \rst{g}}
= \rng{f} \circ \rst{g}
=\rst{\rng{f} \circ g}
= \rst{g} \circ \rng{f}.
$
\end{lemma}
\begin{proof}
\begin{enumerate}[(i)]
\item
That $\rng{f \circ \rst{g}}
= \rng{f} \circ \rst{g}$ is simply axiom RR.3.
\item
That $\rst{\rng{f} \circ g}=\rng{f} \circ \rst{g}$ follows by applying
RR.3 to the occurence of $\rng{f}$ in the left hand side and then by application of R.1.
\item
That $\rng{f \circ \rst{g}}=\rst{g} \circ \rng{f}$ follows by using 
RR.1, then R.2 and then RR.1 again.
\end{enumerate}
\end{proof}
\begin{lemma}
\llabel{fglemma}
 \begin{enumerate}[(i)]
 \item %(i)
 If \fgsinkdiag in \catcw then
 $\rng{f}\circ \rng{g} = \rng{g} \circ \rng{f},$
 \item %(ii)
 If $\sequentialdiag{a}{b}{c}{f}{g}$ in \catcw  then
$\rng{g} \circ \rng{f \circ g}=  \rng{f \circ g}.$
 \item %(iii)
If \fgsinkdiag in \catcw  then
$\rng{f \circ \rng{g}}=\rng{f} \circ \rng{g}.$
 \item %(iv)
 If \fgsourcediag in \catcw  then
$\rst{f \circ \rng{\rst{g} \circ f}} \circ g = \rst{f} \circ g.$
 \commentary{Original (iv) moved to lemma \lref{fgJoinLemma}}
\item %(v)
 If $\sequentialdiag{a}{b}{c}{f}{g}$ in \catcw  
 then if $\rst{g} \leq \rng{f}$ then $\rng{f \circ g} = \rng{g}$.
  \item %(vi)
 If $\sequentialdiag{a}{b}{c}{f}{g}$ in \catcw  then
$\rst{f \circ g} \leq \rst{f}.$
 \item %(vii) 
 If  $\binarysinkContinued{a}{b}{c}{d}{f}{g}{h}$
in \catcw then
 $\rng{ f \circ \rng{g} \circ h } = \rng{ g \circ \rng{f} \circ h } .$
\end{enumerate}
\end{lemma}
\begin{proof}
\begin{enumerate}[(i)]
\item Follows from RR.1 and R.2.
\item 
\begin{align*}
\rng{g} \circ \rng{f \circ g}&=\rng{f \circ g} \circ \rng{g}&&\mbox{by (i)} \\
&=\rng{f \circ g} \circ \rst{\rng{g}}
                    &&\mbox{by lemma \ref{unirangelemma} (\ref{lem:unirange:hatbar}),} \\
&= \rng{f \circ g \circ \rst{\rng{g}}}&&\mbox{by RR.3,} \\
&= \rng{f \circ g} &&\mbox{by lemma \ref{unirangelemma} (\ref{lem:unirange:hatbar}), then RR.2.}
\end{align*}
\item %(iii)
\begin{align*}
\rng{f \circ \rng{g}} 
               &=\rng{f \circ \overline {\rng{g}}} &\mbox{by RR.1}\\
               &=\rng{f} \circ \overline {\rng{g}} &\mbox{by RR.3}\\
               &=\rng{f} \circ \rng{g}             & \mbox{by RR.1}\\
\end{align*}
\item %{iv}
\begin{align*}
\rst{f \circ \rng{\rst{g} \circ f}} \circ g
               &= \rst{f \circ \rng{\rst{g} \circ f}} \circ \rst{g} \circ g
               && \mbox{by R.1,}\\
               &= \rst{\rst{f \circ \rng{\rst{g} \circ f}} \circ \rst{g}} \circ g
               && \mbox{by R.3 then wR.4,}\\
               &= \rst{\rst{g} \circ \rst{f \circ \rng{\rst{g} \circ f}}} \circ g
               && \mbox{by R.2,}\\
               &= \rst{\rst{\rst{\rst{g}} \circ f \circ \rng{\rst{g} \circ f}}} \circ g
               && \mbox{by R.3,}\\
               &= \rst{\rst{g} \circ f \circ \rng{\rst{g} \circ f}} \circ g
               && \mbox{by lemma \ref{unirangelemma} (i),}\\
               &= \rst{\rst{g} \circ f } \circ g
               && \mbox{by RR.2,}\\
               &= \rst{g} \circ \rst{f}  \circ g
               && \mbox{by R.3,}\\
               &=  \rst{f}  \circ \rst{g}  \circ g
               && \mbox{by R.2}\\
               &= \rst{f}  \circ g
               && \mbox{by R.1, as required.}
\end{align*}
\item %(v)
From the inital assumption that
$\rst{g} \leq \rng{f}$ which, by defintion and by  RR.1 and R.2 can be restated as
$\rng{f} \circ \rst{g}  = \rst{g}$,
we are required to prove that
$\rng{f \circ g} = \rng{g}$
which we do as follows:
\begin{align*}
\rng{f \circ g} 
               &= \rng{\rng{f} \circ g}
               && \mbox{by RR.4,} \\
               &= \rng{\rng{f} \circ \rst{g} \circ g}
               && \mbox{by R.1,} \\
               &= \rng{\rst{g} \circ g}
               && \mbox{from the initial assumption restated,} \\
               &= \rng{g}
               && \mbox{by R.1, as required.} 
\end{align*}
\item %(vi)
Need show $\rst{f \circ g} \circ \rst{f} = \rst{f \circ g}$
which we can do as follows
\begin{align*}
\rst{f \circ g} \circ \rst{f} 
                 &= \rst{\rst{f \circ g} \circ f}
                            && \mbox{by R.3,} \\
                 &= \rst{f \circ \rst{g} } 
                            && \mbox{by R.4,} \\
                 &= \rst{f \circ g}
                            && \mbox{by wR.4.}
\end{align*}
\item %vii

\begin{align*}
\rng{ f \circ \rng{g} \circ h } 
            &= \rng{ \rng{f} \circ \rng{g} \circ h }
            && \mbox{by RR.4,}\\
            &= \rng{ \rng{g} \circ \rng{f} \circ h }
            && \mbox{by part (i) of this lemma,}\\
            &= \rng{ g \circ \rng{f} \circ h }
            &&\mbox{by RR.4, as required.}
\end{align*}
\end{enumerate}
\end{proof}


\begin{lemma}
\llabel{tighteningashorteninglemma}
If 
$
% In this array 
% (i) the \nudgeup is so that space is allocated for the
% morphisms looping above the real estate of the array
\begin{array}{c c c c}
\nudgeup{0.75cm}
          & \bOb{y} &         & \bOb{z} \\[0.5cm]
  \bOb{w} &         & \bOb{x} & \\[0.4cm]
\end{array}
\begin{arrows}
\ncarr{w}{y}\alabel{f}
\ncarr{x}{y}\alabel{g}
\ncarr{x}{z}\blabel{h}
\drawMorphismArcing{y}{z}{a}{s}{15}{15}[0.5]
\ncarrrecursive{w}{w}{270}{0}{5.5}\alabel{\overline{i}} 
\ncarrrecursive{z}{z}{90}{180}{8.5}\alabel{\overline{j}}  
\end{arrows} 
$
in a range category \catcw and if
$g \circ s = h$  
then
\begin{enumerate}[(i)]
\item 
\llabel{tighteningshorteningRightToLeft}
 if $\rng{f}=\rng{g}$ then
$
f \circ \rng{\overline{h \circ j} \circ g} \circ s = f \circ s \circ \rst{j},
$
\item 
\llabel{tighteningshorteningLeftToRight}
 if $\rng{f}\leq\rng{g}$ then
$\rng{\overline{g \circ \rng{\rst{i} \circ f}}\circ h} = \rng{\rst{i} \circ f \circ s}$.
\end{enumerate}
\end{lemma}
\begin{proof}
\begin{enumerate}[(i)]
\item

\begin{align*}
f \circ \rng{\overline{h \circ j} \circ g} \circ s
   &= f \circ \rng{\overline{g \circ s \circ j} \circ g} \circ s
                     &&\mbox{substituting $h=g \circ s$, } \\
   &= f \circ \rng{g \circ  \overline{s \circ j} } \circ s
                     && \mbox{by R.4}, \\
   &= f \circ \rng{g} \circ  \overline{s \circ j}  \circ s 
                     && \mbox{by RR.3}, \\
   &= f \circ s \circ \rst{j}
                     && \mbox{because $\rng{f} = \rng{g}$ and by R.4, as required.} 
\end{align*}
\item

\begin{align*}
\rng{\overline{g \circ \rng{\rst{i} \circ f}}\circ h}
     &= \rng{\overline{g \circ \rng{\rst{i} \circ f}}\circ g \circ s}
     && \mbox{substituting $h=g \circ s$,} \\ 
     &= \rng{g \circ \rng{\rst{i} \circ f}\circ s}
     && \mbox{by R.4 and RR.1,} \\ 
     &= \rng{\rng{g} \circ \rng{\rst{i} \circ f}\circ s}
     && \mbox{by RR.4,} \\ 
     &= \rng{\rng{g} \circ \rng{f} \circ \rng{\rst{i} \circ f}\circ s}
     && \mbox{by lemma \ref{fglemma} (ii),} \\ 
     &= \rng{\rng{f} \circ \rng{\rst{i} \circ f}\circ s}
     && \mbox{from the assumption that $\rng{f} \leq \rng{g}$,}\\
     &= \rng{\rng{\rst{i} \circ f}\circ s}
     && \mbox{by lemma \ref{fglemma} (ii),} \\ 
     &= \rng{\rst{i} \circ f \circ s}
     && \mbox{by RR.4, as required.}   
\end{align*}
\end{enumerate}
\end{proof}